\section{Koncentrace a titrace}

\subsection{Koncentrace}

\subsubsection{Molární koncentrace}

Podíl látkového množství rozpuštěné látky a celkového objemu roztoku.

$c = \frac{n}{V} = \frac{m}{M . V}\ \ \ [mol.dm^{-3}]$

Ukázat odměrnou baňku a pipetu, kalibrace EX a IN.

\subsubsection{Molální koncentrace}

Podíl látkového množství rozpuštěné látky a hmotnosti rozpouštědla.

$c_M = \frac{n}{m_S} = \frac{m}{M . m_S}\ \ \ [mol.kg^{-1}]$

\subsubsection{Hmotnostní zlomek}

Podíl hmotnosti složky směsi a celkové hmotnosti směsi. Bezrozměrná veličina.

$w_i = \frac{m_1}{\sum m_i}$

\subsubsection{Molární zlomek}

Podíl látkového množství složky směsi a součtu látkových množství všech složek směsi. Bezrozměrná veličina.

$X_i = \frac{n_1}{\sum n_i}$

\subsection{Titrace}

\zadani{0,1315 g monohydrátu šťavelanu vápenatého bylo titrováno v prostředí kyseliny sírové 0,02 M roztokem \ce{KMnO4}, jehož průměrná spotřeba činila 17,90 cm$^3$. Vypočítejte čistotu šťavelanu.}

\ce{2 KMnO4 + 5 CaC2O4 + 8 H2SO4 -> K2SO4 + 2 MnSO4 + 5 CaSO4 + 10 CO2 + 8 H2O}

\begin{itemize}
	\item n(\ce{KMnO4}) = c.V = 0,02 . 0,01790 = 0,000358 mol
	\item n(\ce{CaC2O4}) = 0,000895 mol
	\item M(\ce{CaC2O4.H2O}) = 146,11 g.mol$^{-1}$
	\item m(\ce{CaC2O4.H2O}) = 0,1308 g
	\item w = $\frac{0,1308}{0,1315} = 0,9945$
\end{itemize}

\odpoved{Čistota šťavelanu je 99,45 \%.}

\clearpage
\zadani{Z navážky 0,6239 g znečištěného \ce{K2CrO4} bylo v odměrné baňce připraveno 100 cm$^3$ roztoku. Z tohoto roztoku bylo odpipetováno 20,0 cm$^3$, přidána destilovaná voda, KI, zředěná HCl a roztok škrobu. Vyloučený jod byl titrován 0,1 M roztokem \ce{Na2S2O3}, jehož spotřeba na titraci činila 18,90 cm$^3$. Vypočítejte procentuální obsah \ce{K2CrO4} v navážce.}

Jde o nepřímou titraci:

\ce{2 K2CrO4 + 6 KI + 16 HCl -> 2 CrCl3 + 10 KCl + 3 I2 + 8 H2O}

\ce{2 Na2S2O3 + I2 -> Na2S4O6 + 2 NaI}

\ce{6 Na2S2O3 + 3 I2 -> 3 Na2S4O6 + 6 NaI} (po vynásobení třemi, kvůli jódu)

\begin{itemize}
	\item M(\ce{K2CrO4}) = 194,19 g.mol$^{-1}$
	\item n(\ce{Na2S2O3}) = c.V = 0,1 . 0,01890 = 0,00189 mol
	\item n/6 = 0,000315
	\item n(\ce{K2CrO4}) = 0,000315 . 2 = 0,00063 mol
	\item m(\ce{K2CrO4}) = 0,00063 . 194,19 = 0,1223 g
	\item Zřeďovací faktor: m = 0,1223 . 5 = \textbf{0,6117 g} \ce{K2CrO4}
	\item w = $\frac{0,6117}{0,6239} = 0,9804$
\end{itemize}

\odpoved{Obsah chromanu byl 98,04 \%.}